\section{Replicação do PGDCS contra a fonte}

O objetivo~4 concluiu, dentro do protocolo deste projeto, que votar as medidas de
complexidade não muda \On{1}, \mvr{} nem \dfe{} de forma sustentada, e o objetivo~8
traduziu isso numa recomendação: o PGDCS completo não paga o custo que cobra. A fonte do método relata o oposto --- ganhos de $4$ a $6$ pontos percentuais
sobre Bagging na Tabela~A.10 de \textsc{Monteiro et al.}~(2023). Duas afirmações
incompatíveis sobre o mesmo método exigem reconciliação, e ela só é possível
replicando o protocolo do artigo, não o deste relatório.

Esta seção reporta essa replicação. Ela usa árvore de decisão (não Perceptron) e
duas configurações de protocolo registradas na ADR~0020, executadas em
\texttt{7c5ec6d}:

\begin{description}
\item[Perfil do artigo.] \texttt{monteiro-2023-\hspace{0pt}dt-\hspace{0pt}compatible-v1}: $20$
\emph{holdouts} estratificados independentes $50/25/25$, semente-raiz $86$,
validação \emph{coincidente} com o DSEL, $M=100$, $20$ gerações. Run
\texttt{2026-08-27\hspace{0pt}T19\hspace{0pt}-47-09}: $1620$ linhas, $27$ bases, $3$h$14$.
\item[Perfil do POS.] \texttt{pos-oracle-v2}: $10$-fold estratificado com
partição interna treino / validação / DSEL \emph{disjunta}, semente $42$, mesmo
indutor e mesmo $M$. Run \texttt{2026-08-27\hspace{0pt}T23\hspace{0pt}-01-32}: $840$ linhas, $28$ bases,
$5$h$38$.
\end{description}

\subsection{O gap, e o que o produz}

Contra a Tabela~A.10, os três tratamentos se comportam de modo qualitativamente
diferente:

\begin{table}[htbp]
\centering\small
\caption{Diferença contra a Tabela~A.10 no perfil compatível com o artigo, $27$
bases. O que decide a questão é o \emph{sinal}, não a magnitude.}
\begin{tabular}{lrrr}
\toprule
Tratamento & erro médio $|\Delta|$ & média com sinal & negativos \\
\midrule
PGDCS\_A & $6{,}19$ pp & $\mathbf{-6{,}19}$ pp & $\mathbf{27/27}$ \\
PGDCS\_D & $5{,}25$ pp & $\mathbf{-5{,}25}$ pp & $\mathbf{27/27}$ \\
Bagging  & $0{,}75$ pp & $-0{,}16$ pp & $14/27$ \\
\bottomrule
\end{tabular}
\end{table}

O Bagging --- que compartilha \emph{split}, semente, indutor, $M$ e protocolo de
avaliação com os outros dois --- erra $0{,}75$ pp em módulo, com média com sinal de
$-0{,}16$ pp e $14$ negativos em $27$: ruído simétrico, a assinatura de uma
reimplementação fiel. O PGDCS erra sempre para baixo, em $54$ comparações
consecutivas, sem uma exceção. Um defeito de implementação dispersaria nos dois
sentidos; um viés sistemático, não.

\paragraph{A causa é viés de seleção na medida publicada.} O Algoritmo~1 do artigo
escolhe a geração final por $\upsilon \leftarrow \text{acurácia}(E^{(t+1)}, V)$, e a
Tabela~A.10 reporta essa mesma quantidade --- o máximo sobre $20$ gerações, medido
no \emph{mesmo} conjunto de $25\%$ que fez a escolha. O número do Bagging, que não
passa por seleção alguma, é uma medida honesta em teste. Gravando as trajetórias
geração a geração em Wine, Heart, Blood e Sonar ($20$ replicações cada), a
distinção fica direta: chamando \texttt{VAL@sel} a acurácia da geração escolhida no
conjunto de validação e \texttt{TEST@sel} a da mesma geração em teste, o publicado
dista $\approx1{,}7$ pp de \texttt{VAL@sel} e $\approx6{,}3$ pp de
\texttt{TEST@sel}.

Três verificações independentes fecham o diagnóstico:

\begin{itemize}
\item \textbf{A trajetória tem a forma de sobreajuste ao conjunto de seleção.} A
\mvr{} de validação sobe monotonicamente até a geração $20$; a acurácia de teste
sobe até a geração $\approx4$--$8$ e depois cai (Wine: validação $95{,}0 \to
99{,}8$, teste $95{,}3 \to 95{,}0$).

\item \textbf{O teto por replicação não alcança o publicado.} Tomando, em cada
replicação, a \emph{melhor} geração medida em teste --- um oracle sobre gerações,
inatingível por qualquer critério de parada --- Wine chega a $98{,}0$/$97{,}4$,
ainda abaixo dos $99{,}3$ publicados. Nenhuma política de seleção fecharia o gap.

\item \textbf{O que não depende do critério de saída reproduz.} \On{1} bate a
Tabela~7 com erro médio de $0{,}44$ pp e a geração escolhida bate a Tabela~6 com
$3{,}19$ gerações de erro. O Oracle reproduz porque não passa pela seleção.
\end{itemize}

\subsection{O tamanho honesto do efeito}

Com os dois perfis executados, dá para medir o que o PGDCS ganha do Bagging quando
as duas pontas são medidas em teste --- o contraste que a Tabela~A.10 não faz:

\begin{table}[htbp]
\centering\small
\caption{PGDCS\_A menos Bagging, acurácia \mvr{} em teste, pareado por base
(Wilcoxon). O ganho existe, é uma ordem de grandeza menor que o publicado, e
depende do tamanho do conjunto de validação.}
\begin{tabular}{llrrr}
\toprule
Perfil & Validação & $\Delta$ médio & Vence em & $p$ \\
\midrule
\texttt{monteiro-2023-dt} & $25\%$, $=$ DSEL & $\mathbf{+0{,}39}$ pp & $22/27$ & $\mathbf{0{,}0030}$ \\
\texttt{pos-oracle-v2} & $10\%$, disjunta & $-0{,}13$ pp & $12/27$ & $0{,}75$ \\
\bottomrule
\end{tabular}
\end{table}

No perfil compatível com o artigo o ganho existe e é consistente, mas vale
$\approx0{,}4$ pp contra os $4$--$6$ pp anunciados. No perfil deste projeto, cuja
validação é menor e disjunta do DSEL, ele desaparece: $12$ vitórias em $27$ é
sorteio ($p=0{,}75$). A seleção de geração precisa de um conjunto de validação
grande para não escolher ruído, e o que entrega, quando entrega, é marginal.

\textbf{Isto confirma a conclusão do objetivo~4 por um caminho independente.} Lá, o
PGDCS completo era indistinguível dos bags aleatórios em pools de Perceptron sob o
protocolo do POS; aqui, com árvore de decisão e sob o protocolo do próprio artigo,
o ganho publicado se decompõe em viés de seleção mais $0{,}4$ pp reais. As duas
medições não estavam em conflito com a fonte --- a fonte é que compara uma
quantidade otimizada com uma não otimizada.

\paragraph{Nenhuma mudança de protocolo foi feita para fechar o gap}, porque
fechá-lo exigiria reportar a métrica enviesada. A coluna comparável à Tabela~A.10
(\mvr{} de validação da geração vencedora) fica gravada no manifesto de fold como
\texttt{pgdcs\_best\_score}; \texttt{mv\_acc} continua sendo acurácia em teste.

\subsection{As medidas exatas do ECoL não mudam o quadro}

O backend rápido de complexidade não reproduz quatro das onze medidas do ECoL
0.3.0, a versão que o artigo usou: T1, N1, F4 e LSC (ADR~0023). Para testar se
isso explica parte do gap, o perfil \texttt{monteiro-2023-dt-compatible-v2}
repete o v1 com as portas exatas dessas quatro medidas na votação e no fitness
do AG (ADR~0024); \emph{splits}, sementes e o resto do protocolo são os mesmos.

\begin{table}[htbp]
\centering\small
\caption{Perfil compatível com o artigo com as medidas do backend rápido (v1) e
com as do ECoL 0.3.0 (v2), $27$ bases, $20$ replicações.}
\begin{tabular}{lrr}
\toprule
 & v1 & v2 \\
\midrule
Vencedor overlapping da Tabela~4 acertado & $12/27$ & $13/27$ \\
Vencedor neighborhood da Tabela~4 acertado & $7/27$ & $8/27$ \\
\mvr{} PGDCS\_A $-$ Tabela~A.10 & $-6{,}19$ pp & $-6{,}39$ pp \\
\mvr{} PGDCS\_D $-$ Tabela~A.10 & $-5{,}25$ pp & $-5{,}29$ pp \\
PGDCS\_A $-$ Bagging (teste) & $+0{,}39$ pp, $p=0{,}003$ & $+0{,}20$ pp, $p=0{,}058$ \\
\bottomrule
\end{tabular}
\end{table}

A votação muda em $18$ das $27$ bases, mas a concordância com a Tabela~4 quase
não se move, e o PGDCS continua $\approx6$ pp abaixo do publicado em todas as
bases. As definições das medidas ficam descartadas como causa do gap. O ganho do
PGDCS sobre o Bagging cai pela metade e deixa de ser significativo a $5\%$; a
diferença entre v2 e v1 por base não é significativa ($-0{,}20$ pp, $p=0{,}33$),
mas o $+0{,}4$ pp da tabela anterior depende da definição das medidas.

\subsection{Duas correções de implementação encontradas no caminho}

A replicação expôs dois defeitos no PGDCS deste projeto, ambos corrigidos por TDD e
sem efeito sobre o diagnóstico acima:

\begin{description}
\item[Votação das medidas (Eq.~10--11).] O voto usava a \emph{média acumulada} dos
$\sigma_j^t$ em vez do $\sigma_j^t$ da própria iteração, o que concentrava $9$ ou
$10$ dos $10$ votos de uma família numa única medida --- contra o espalhamento que
as Tabelas~3 e~4 do artigo mostram. Corrigido, o acerto contra a Tabela~4 sobe de
$1/5$ bases para $4/8$ medidas nas quatro primeiras bases. O efeito na acurácia é
de $+0{,}6$ a $+1{,}3$ pp, só no PGDCS: a escolha das medidas não é a causa do gap.

\item[A fração $u$.] §5.1 e §5.2 do artigo descrevem passos diferentes na mesma
frase. Adotado $u=0{,}4$ para as sub-amostras da votação e $0{,}5$ para os
\emph{bags} do AG; antes ambos valiam $0{,}5$.
\end{description}
